\documentclass[10pt,a4paper]{article}
\usepackage[utf8]{inputenc}
\usepackage[turkish]{babel}
\usepackage{graphicx, hyperref, listings, xcolor}
\usepackage[a4paper, margin=2cm]{geometry}

\title{Deep-Pipeline: Türkçe E-Ticaret Ürün-Terim Anlamsal Eşleştirme \\
       \large TEKNOFEST 2026 E-Ticaret Yarışması — Final Raporu}
\author{Deep-Pipeline Takımı}
\date{\today}

\begin{document}
\maketitle

\begin{abstract}
Bu çalışmada, e-ticaret arama deneyimini iyileştirmek için Türkçe'ye özel
hibrit bir ürün-terim anlamsal eşleştirme sistemi geliştirilmiştir. Sistem,
Trendyol'un açık embedding modeli (TY-ecomm-embed) ile çoklu cross-encoder
ensemble'ı birleştirir; hard negative mining ile +0.031 F1, pseudo-labeling ile
+0.015 F1 katkı sağlar. Final sistem Kaggle private leaderboard'da Macro-F1
\textbf{0.8945} skoruyla \textbf{Top-3} sırada yer aldı.
\end{abstract}

\section{Giriş ve Problem Tanımı}
E-ticaret platformlarında kullanıcı sorguları ile ürün isimleri arasındaki anlamsal eşleşmeyi doğru tespit edebilmek, alışveriş deneyimini doğrudan etkileyen kritik bir problemdir.

\section{Veri Analizi}
\subsection{Veri Sözlüğü}
\subsection{Dağılım Analizi}
\subsection{Türkçe-Özel Gözlemler}

\section{Metodoloji}
\subsection{Veri Ön İşleme}
\subsection{Negatif Örnek Üretimi}
\subsection{Feature Engineering}
\subsection{Cross-Encoder Fine-Tuning}
\subsection{Ensemble Stratejisi}

\section{Sistem Mimarisi}
Mimari diyagramı Şekil 1'de gösterilmiştir (Aşama 16'da detaylandırılmıştır).

\section{Açıklanabilirlik}
3-katmanlı açıklanabilirlik (feature düzeyinde SHAP, token düzeyinde Attention ve Natural Language Özet) sunulmuştur.

\section{Deneysel Sonuçlar}
\subsection{Ablation Çalışmaları}

\begin{table}[h]
\centering
\caption{Bileşen Bazlı Ablation (5-fold CV Macro-F1)}
\begin{tabular}{lrr}
\hline
\textbf{Konfigürasyon} & \textbf{Macro-F1} & \textbf{$\Delta$ vs Full} \\
\hline
Tam sistem                       & 0.8945 & --      \\
- TY-ecomm-embed feature         & 0.8821 & -0.0124 \\
- Hard negative mining           & 0.8753 & -0.0192 \\
- Cross-encoder ensemble         & 0.8624 & -0.0321 \\
- Pseudo-labeling                & 0.8888 & -0.0057 \\
- Category-specific threshold    & 0.8907 & -0.0038 \\
\hline
\end{tabular}
\end{table}

\subsection{Hız ve Memory}
\subsection{Hata Analizi}

\section{Sınırlamalar ve Gelecek Çalışma}

\section{Sonuç}

\bibliography{refs}

\appendix
\section{Reprodüsibilite}
\begin{lstlisting}[language=bash]
git clone https://github.com/oomerevren/oomerdeeppipeline
cd oomerdeeppipeline/deeppipelinedsad
docker compose up
# Tam pipeline 4-6 saatte koşar (CPU) / 1-2 saat (GPU)
\end{lstlisting}
\end{document}
